\uuid{TQuJ}
\titre{Pour tout/il existe}
\auteur{bodin ; hanani ; mzari}
\organisation{exo7}
\licence{CC BY-NC-SA 4.0}
\datecreate{2020-01-20}
\niveau{L1}
\module{Algèbre}
\chapitre{Logique, ensemble, raisonnement}
\sousChapitre{Logique}
\difficulte{3}
\theme{Logique, pour tout, il existe}
\exo7qcmid{lille-1.logique.016}

\contenu{
\begin{question}
\qclassification{100.01}
\qtags{motcle=pour tout/il existe}
Soit $f : \Rr \to \Rr$ la fonction définie par $f(x)=\exp(x)-1$.
Quelles sont les assertions vraies ?
\begin{answers}
    \good{$\forall x,x' \in \Rr  \qquad x \neq x' \implies f(x) \neq f(x')$}
    \good{$\forall x,x' \in \Rr  \qquad x \neq x' \Longleftarrow f(x) \neq f(x')$}
    \bad{$\forall x,x' \in \Rr  \qquad x \neq x' \implies (\exists y \in \Rr \quad f(x) < y < f(x'))$}
    \good{$\forall x,x' \in \Rr  \qquad  f(x)\times f(x') < 0 \implies x\times x' < 0$}
\end{answers}
\begin{explanations}
Dessiner le graphe de $f$ pour mieux comprendre !
Même si $f(x) \neq f(x')$ cela ne veut pas dire que $f(x) < f(x')$, l'inégalité pourrait être dans l'autre sens.
\end{explanations}
\end{question}
}
